% ==============================================================================
% Section: Cap-Incidence Across Regimes
% Status: NEW (2026-04-26 — Phase A of cap-incidence pivot)
% ==============================================================================

\section{Cap-Incidence Across Regimes}
\label{sec:cap_incidence}

This section documents which administrative constraint was binding in each
of PJM's seven completed Base Residual Auctions: the TPS mitigation cap
based on Net ACR, the design-level VRR Point~(a) cap, or the Shapiro
settlement cap that took effect for the 2026/27 and 2027/28 delivery
years. The cap-incidence panel reframes the dramatic price changes
observed across the sample period as a regime change in the binding
administrative constraint rather than a behavioral shift in seller
strategy.

\begin{table}[H]
  \centering
  \caption{Cap-Incidence Panel: Which Administrative Constraint Was Binding,
    PJM Base Residual Auctions 2021/22--2027/28 (RTO Level)}
  \label{tab:cap_incidence}
  \small
  \begin{tabular}{lcccccrcl}
    \toprule
             & VRR     & Clearing & Net      & VRR     & Settle.   & Reserve   & TPS     & Binding   \\
    Year     & Design  & price    & ACR      & Pt~(a)  & cap       & excess IRM (MW) & mitig.\ & cap       \\
    \midrule
    2021/22 & Old & \$140 & \$150 & \$482 & --- & --- & --- & None \\
    2022/23 & Old & \$50 & \$150 & \$391 & --- & --- & 0.0\% & None \\
    2023/24 & Old & \$34 & \$150 & \$412 & --- & +5,687 & 7.3\% & TPS active \\
    2024/25 & Old & \$29 & \$150 & \$440 & --- & +3,278 & 2.2\% & TPS active \\
    2025/26 & Old & \$270 & \$150 & \$452 & --- & -205 & 5.5\% & TPS active \\
    2026/27 & New & \$329 & \$150 & \$329 & \$325 & -492 & 6.3\% & Settlement \\
    2027/28 & New & \$333 & \$150 & \$333 & \$325 & -6,516 & 1.0\% & Settlement \\
    \bottomrule
  \end{tabular}
  \vspace{3pt}

  \begin{minipage}{0.96\textwidth}
    \footnotesize \textit{Notes:} Prices in nominal \$/MW-day UCAP.
    Net ACR is the IMM-published sector-average Avoidable Cost Rate
    (\$149.32, rounded to \$150) used as the default mitigation offer cap
    for resources subject to the Three Pivotal Supplier (TPS) test.
    VRR Pt~(a) is the design-level price cap on the unmodified Variable
    Resource Requirement curve. Settlement cap is the Shapiro/PJM
    \$325/MW-day UCAP cap that took effect for the 2026/27 and 2027/28
    delivery years. Reserve excess IRM (MW) is committed UCAP minus
    IRM-implied required reserves (negative = shortfall). TPS mitigation
    percentage is the share of generation resources that submitted Capacity
    Performance offers and had unit-specific Market Seller Offer Caps
    calculated by the IMM. Binding cap classification: ``Settlement'' =
    cleared within 2\% of the new-design Point~(a), which under the
    Shapiro/PJM agreement embodies the settlement cap; ``TPS active'' =
    TPS mitigation calculated for $\geq 1$\% of resources with clearing
    more than 5\% below Point~(a), indicating mitigation operative for
    some marginal resources but not pinning headline clearing; ``None'' =
    clearing well below all administrative caps with TPS mitigation
    calculated for $<1$\% of resources.
    Sources: PJM BRA results (clearing prices, cleared MW, reliability
    requirement); IMM State of the Market Reports 2021--2025, Section 5
    Capacity (TPS mitigation counts, reserve-margin panel).
    The 2021/22 BRA (held May 2018) predates locally-available IMM SoMs,
    so TPS mitigation and reserve-excess data are reported as ``---''.
  \end{minipage}
\end{table}


Three regimes are visible in Table~\ref{tab:cap_incidence}, distinguished
by the joint pattern of clearing price, capacity margin, and IMM-applied
mitigation activity.

The first regime, covering 2021/22 through 2024/25, is one of comfortable
supply margins and cost-anchored clearing. The IMM-published reserve
margin in excess of the Installed Reserve Margin (IRM) was strongly
positive in 2023/24 and 2024/25 (+5{,}687~MW and +3{,}278~MW respectively;
data for 2021/22 and 2022/23 are not in the IMM panel I assembled), and
clearing prices ranged from \$29 to \$140/MW-day, well below both the
sector-average Net ACR and the design-level VRR Point~(a) cap of
\$391--\$482/MW-day. TPS mitigation was applied to a small share of
generation resources (0.0\% in 2022/23, 7.3\% in 2023/24, 2.2\% in
2024/25), but with the marginal cleared offer well below any active cap,
mitigation operative for a minority of supra-marginal resources did not
pin the clearing outcome. The classification ``TPS active'' for 2023/24
and 2024/25 records the IMM mitigation activity without claiming that the
mitigation cap determined the clearing price; ``None'' for 2021/22 and
2022/23 records a slack market in which mitigation was either inactive
or not reported in locally available IMM data.

The second regime, in 2025/26, is a transition. The reserve margin in
excess of IRM turned slightly negative for the first time
(\textminus205~MW), the clearing price jumped from \$29 in the
preceding auction to \$270/MW-day, and TPS mitigation rose to 5.5\% of
resources. Yet \$270 remains far below the design Point~(a) of
\$452/MW-day, so no administrative cap was binding the clearing.
2025/26 is the auction in which the underlying supply--demand balance
tightened materially without any cap yet binding the headline outcome.

The third regime, comprising 2026/27 and 2027/28, is one of
cap-binding clearing under negative reserve excess. Both auctions
cleared at the new-design VRR Point~(a) value, which under the
Shapiro/PJM settlement embodies the agreed maximum price (the nominal
\$325/MW-day target translates to \$329.17 and \$333.44 in 2026/27 and
2027/28 respectively after year-specific accreditation and Net~CONE
adjustments). The reserve margin in excess of IRM was \textminus493~MW
in 2026/27 and \textminus6{,}517~MW in 2027/28, with the IMM declaring
both auctions ``not competitive''
\citep{MonitoringAnalytics2025_sotm}.\footnote{The IMM SoM 2025
reports the 2026/27 reserve shortfall in two forms: the panel-consistent
``Reserve cleared in excess of IRM UCAP'' figure of \textminus493~MW
used here, and a parallel calculation reported as ``208.7~MW short of
the required reserve level of 21{,}561.9~MW'' computed against a
different reserve baseline. Both characterize a marginal shortfall.}
The drop in TPS mitigation application from 6.3\% in 2026/27 to 1.0\%
in 2027/28 is consistent with the settlement cap moving sufficiently
close to the underlying offer distribution that fewer individual offers
exceeded their unit-specific Net ACR caps.

The cap-incidence panel offers a direct answer to a recurring question
about the post-settlement auctions: did seller behavior change between
the slack and tight regimes, or did the binding administrative
constraint change while the underlying competitive process remained
similar? The data favor the latter interpretation. Across the sample,
the SFE machinery developed in Section~\ref{sec:model} predicts
strategic markups that are far above observed slack-year clearings and
mechanically equal to the cap in tight years. The honest reading is
that the SFE identifies the at-cap versus interior \emph{regime},
conditional on supply margin and concentration, but does not
independently identify the dollar clearing price in either regime.
What identifies clearing in the data is the lowest binding
administrative constraint: cost in slack years, the settlement cap in
tight years. The apparent dramatic price change between regimes is
therefore at least in part a regime change in which constraint binds,
not a structural shift in strategic conduct.

This reading does not contradict the IMM's substantive assessment. The
IMM attributes the underlying scarcity in 2026/27 and 2027/28 to data
center load growth rather than to any change in seller behavior or to
the settlement itself \citep{MonitoringAnalytics2025_sotm}. The
contribution of the cap-incidence panel is to make that interpretation
visible at a glance and to anchor the more detailed counterfactual
discussion in Section~\ref{sec:21billion} in observable cap-incidence
facts rather than in model-implied prices that are mechanically tied to
whichever cap is operative.
